\section{Reading Comprehension Questions}

\noindent
{\em From Section~\ref{sec:multSum}}

\begin{requ}
For each of the following, make up an example that requires the given rule to solve.
Your example should not just rehash an example from the book.
Bonus points if it is computer science related.
For each, also give and explain the solution.
\begin{lenum}
\item The sum rule
\item The product rule
\item Both the sum rule and product rule
\end{lenum}
\begin{requanswer}
Answers will vary, but here are a few examples.
(a) I want to adopt a dog. At the pet shop, there are four golden retrievers, two poodles and 5 corgis. 
How many choices do I have if I plan to adopt one dog? 
You can choose either a golden retriever (4 options), a poodle (2 options), or a corgi (5 options).
So the total number of choices is $4 + 2 + 5 = 11$.

(b) I go to a restaurant and see that there are 15 dessert choices, 
10 main courses, and 2 appetizers. How many choices do I have if I want to order one of each?
You have 15 choices for desserts, and independent of that you have 10 choices for a main course,
and independent of both of those you can choose one of two appetizers. 
So you have $15 * 10 * 2 = 300$ choices.

(c) Your password can be from 4 to 8 lower case alphabetic characters. How many possibilities are there?
There are 26 possible characters. If you use 4 characters, there are 
$26\times 26\times 26\times 26 = 26^4$ possible passwords.
Similarly, with 5 characters there are $26^5$ possible passwords.
Likewise, for 6, 7, and 8 characters, there are $26^6$, $26^7$, and $26^8$ possible passwords.
Since you have to choose one of these lengths, the total number of possible passwords is
$26^4+26^5+26^6+26^7+26^8$. 
\end{requanswer}
\end{requ}

\noindent
{\em From Section~\ref{sec:pigeon}}

\begin{requ}
If there are 12 objects in 10 boxes, does the pigeonhole principle allow you to conclude that
one box has at least 3 objects? Or that two boxes have at least two items? Or that every box has at least one item? 
What is the most precise thing that you can conclude from it?
\begin{requanswer}
You cannot conclude that a box has at least 3 objects, that two boxes have at least two items,
or that every box has at least one item. 
(For each of these you can come up with a distribution of objects in boxes that does not fit the description.)
The most you can conclude is that at least one box has at least 2 objects.
\end{requanswer}
\end{requ}

\begin{requ}
Assume 30 balls are placed in 7 bins. Give several different possibilities for how many balls are in each bin,
trying to make the examples as different from each other as possible. What is true of all of your examples,
as predicted by the generalized pigeonhole principle? 
\begin{requanswer}
You might have all 30 balls in one bin.
You might have 15 balls in one bin, and 15 in a second bin.
You might have 1 balls in each of the first 6 bins and 24 in the 7th bin.
You might have 4 balls in each of the first 6 bins and 6 balls in the 7th bin.
You might have 4 balls in each of the first 5 bins, and 5 balls in each of the final two bins.
Notice that in all of these examples, there is at least one box which has at least $\lceil 30/7\rceil =  5$ balls
as guaranteed by the generalized pigeonhole principle.
\end{requanswer}
\end{requ}

\begin{requ}
I have 21 disc golf discs, including putters, approach discs, fairway drivers, and distance drivers.
Tell me everything you can say for certain about how many of each type of disc I have.
\begin{requanswer}
There are 4 types of discs. Therefore by the generalized pigeonhole principle, I know that I have 
at least $\lceil 21/4\rceil =  6$ discs of at least one type. But that is all I can say. 
I do not know which type I have at least 6 of. For instance, it is possible all 21 are putter, or
all 21 are distance drivers, for instance, so I do not even know if I have a single disc of any type.
\end{requanswer}
\end{requ}

\begin{requ}
Twelve people each pick a number from 1 to 1000. 
Prove that at least two of them picked numbers that have the same number of 1s in their binary representation
or show why it is possible that this is not the case (i.e. give a counterexample).
\begin{requanswer}
The numbers between 1 and 1000 can all be represented with 10 bits. 
A number between 1 and 1000 can have between 1 and 10 bits that are 1s.
So the 12 numbers can each be placed in one of 10 ``bins'' based on how many 
bits they have in their binary representation. Since we are placing 12 numbers in 10 bins,
at least one bin has at least 2 numbers. In other words, two of the numbers have the same number of 1s
in their binary representation.
(None of the numbers can actually have 10 1s because the number with 10 1s is 1023 which is larger than 1000.
So technically there are only 9 bins. But the argument is the same either way.
However, if I said that ten people picked numbers, then we would need to take this into account to solve the problem correctly.)
\end{requanswer}
\end{requ}

\noindent
{\em From Section~\ref{sec:permComb}}

\begin{requ}
What is the difference between a permutation and a combination?
\begin{requanswer}
Permutations are ordered and combinations are not.
More specifically, a permutation is a reordering of objects, 
whereas a combination is a selection of objects. They are basically completely different things.
\end{requanswer}
\end{requ}

\begin{requ}
How many are there of each of the following (where {\em digit} means decimal digit).
\begin{lenum}
\item Three-digit numbers
\item Three-digit numbers with no repeated digits
\item Sets consisting of three digits. (e.g. $\{4, 0, 5\}$)
\item Lists consisting of three digits. (e.g. $[4, 2, 3]$)
\end{lenum}
\begin{requanswer}
(a) $10\cdot 10 \cdot 10 = 10^3=1000$ (assuming you regard $0=000$, $12=012$, etc. as 3 digit numbers).
(b) $10\cdot 9 \cdot 8 = 720$.
(c) Because sets cannot have repeats, there are $\binom{10}{3}=\frac{10\cdot 9 \cdot 8}{3\cdot 2 \cdot 1} = 120$.
{\em Notice that there are 6 times as many three-digit numbers with no repeated digits than sets of three digits because
each set with three digits leads to 6 different three-digit numbers (e.g. the set $\{1, 2, 3\}$ is the same as the set 
$\{3,1,2\}$, for instance, but the numbers 123, 132, 213, 231, 312, 321 are all different.)}
(d) Because lists can have repeats, there are $10\cdot 10\cdot 10=10^3=1000$.
\end{requanswer}
\end{requ}

\begin{requ}
(a) How many permutations are there of the set $\{8, 6, 7, 5, 3, 0, 9\}$?
(b) How many of these permutations begin with 8 and end with 9?
\begin{requanswer}
(a) $7!$.
(b) $5!$.
\end{requanswer}
\end{requ}

\begin{requ}
You know somebody's PIN number uses the digits 3, 3, 6, 6, 8, but you do not know the order.
How many possible PIN numbers have these digits?
\begin{requanswer}
Using Theorem~\ref{thm:genPerms}, there are $\frac{5!}{2!\cdot 2!\cdot 1!} = \frac{120}{4}=30$.
\end{requanswer}
\end{requ}

\begin{requ}
You need to choose a team of 45 people out of a possible 50 people. 
What is probably a much easier way of thinking about this problem?
\begin{requanswer}
Choose 5 people who are {\em not} on the team.
\end{requanswer}
\end{requ}

\begin{requ}
Compute $\binom{25}{22}$ by hand. (Hint: Be smart!)
\begin{requanswer}
$\binom{25}{22} = \binom{25}{3}=\frac{25\cdot 24\cdot 23}{3\cdot 2\cdot 1} 
= \frac{25\cdot (6\cdot 4)\cdot 23}{6} = 25\cdot 4\cdot 23 = 2300$.
\end{requanswer}
\end{requ}

\begin{requ}
The board of directors for the Holland Running Club has 11 members.
The executive board is a subcommittee of the board of directors consisting of 4 members (from the 11).
The executive board consists of the president, vice-president, treasurer, and secretary.
\begin{lenum}
\item 
How many different possibilities are there for the executive board if we do not care which office they hold?
\item 
If we are given the four members of the executive board, how many ways are there of assigning the offices?
\item
How many different possible executive boards are there (choosing from the whole board) 
if we {\em do} care about who is in which office?
\end{lenum}
\begin{requanswer}
(a) $\binom{11}{4}$.
(b) $4!$.
(c) We can choose the 4 members ($\binom{11}{4}$ ways) and then place them into the offices (i.e. order them, so $4!$ ways)
for a total of $\binom{11}{4}\cdot 4! = 11\cdot 10\cdot 9 \cdot 8 = 7920$ ways.
Alternatively, we have 11 choices for president, and once the president has been decided there are not 10 choices for
vice-president, then 9 for treasurer, and finally 8 for secretary, for a total of  $11\cdot 10\cdot 9 \cdot 8 = 7920$ ways
of choosing.
\end{requanswer}
\end{requ}

\begin{requ}
What are two important factors that influence how you go about counting things 
(i.e. help you determine which technique you will use)?
\begin{requanswer}
Whether or not order matters, whether or not there are repetitions, whether or not objects are distinguishable or not. 
\end{requanswer}
\end{requ}

\noindent
{\em From Section~\ref{chap:bino_mio}}

\begin{requ}
\begin{lenum}
\item Simplify $\dsum_{k=0}^{n} \binom{n}{k}10^{k}$.
\item Compute the previous sum for $n=0,1,2,3,4,5$. (Hint: Use your solution to part (a)!)
\item Attempt to make a connection between this question and Pascal's Triangle.
It may be a bit subtle, but it is kind of neat if you see it.
\end{lenum}
\begin{requanswer}
(a) $\dsum_{k=0}^{n} \binom{n}{k}10^{k}= \dsum_{k=0}^{n} \binom{n}{k}10^{k}1^{n-k} = (10+1)^n=11^n$.
(b) $11^0=1, 11^1=11, 11^2=121, 11^3=1331, 11^4=14641, 11^5=161051$.
(c) The rows of Pascal's triangle look kind of like the powers of 11. When the numbers in the triangle are longer than 1 digit,
you have to actually line them up and add them to get the correct result. But the connection is more clear when you
consider the formula for the Binomial Theorem and think about what it says about a row of the triangle.
\end{requanswer}
\end{requ}

\begin{requ}
Use the Binomial Theorem to expand $(2x-3y)^5$, simplifying as much as possible.
\begin{requanswer}
Plugging in $2x$ and $-3y$ into the formula, we obtain
\begin{eqnarray*}
(2x-3y)^5 &=&
\binom{5}{0}(2x)^0(-3y)^5+
\binom{5}{1}(2x)^1(-3y)^4+
\binom{5}{2}(2x)^2(-3y)^3+
\binom{5}{3}(2x)^3(-3y)^2\\
&&+\binom{5}{4}(2x)^4(-3y)^1+\binom{5}{5}(2x)^5(-3y)^0\\
&=& -243y^5 + 810x\, y^4 - 1080x^2y^4+720x^3y^2-  240 x^4y+32x^5.
\end{eqnarray*} 

\noindent
Notice that the negative sign goes inside the parentheses so that it is included in the powers, and that the constants 
are also inside the parentheses so they are included in the powers.
\end{requanswer}
\end{requ}

\begin{requ}
Use the Binomial Theorem to prove that $\dsum_{k=0}^{n}\binom{n}{k}=2^n$.
\begin{requanswer}
Notice that $\dsum_{k=0}^{n}\binom{n}{k}= \dsum_{k=0}^{n}\binom{n}{k}1^k1^{n-k}=(1+1)^n=2^n$,
where the second-to-last step uses the Binomial Theorem.
\end{requanswer}
\end{requ}

\noindent
{\em From Section~\ref{sec:inclExcl}}

\begin{requ}
A rather large family has 12 children, all who attended college.
6 of them took a math class, 5 took a computer science class,
and  4 took neither a math class or a computer science class.
How many took both a math and a computer science class?
\begin{requanswer}
Let $M$ be the set of children who took math and $C$
those who took computer science. 
Notice that $12-4 = 8$ children took either math or computer science, 
Thus, 
$|M|=6$, $|C|=5$, and $|M\cup C|=8$. Therefore, it
must be that $|M\cap C|=6+5-8=3$ children took both.
\end{requanswer}
\end{requ}

\begin{requ}
In a class of 20 students, 7 show up late and 4 sleep during class. 
How many students do neither? Give both a minimum and maximum since there is not
enough information to know for sure.
\begin{requanswer}
It is possible that the 4 students who sleep were also late, in which case 13 students did neither.
On the other extreme, the 4 students who slept are all different than the 7 who came late. 
In that case there are 9 students who did neither.
So at least 9 and at most 13 students did neither.
\end{requanswer}
\end{requ}

\begin{requ}
You want to use Inclusion-Exclusion on 3 sets
(e.g. your goal is to compute how many things are in the union of 3 sets).  
You are given 6 pieces of information. 
Is that enough to solve the problem? Explain.
\begin{requanswer}
No. There are 7 things on the right side of the equation,
so if you only have 6 pieces of information you cannot fully solve the problem.
\end{requanswer}
\end{requ}

\begin{requ}
Given Inclusion-Exclusion on two and three sets, can you generalize it to four sets?
Thus, given sets $A$, $B$, $C$, and $D$, what is a formula for $|A\cup B\cup C\cup D|$?
\begin{requanswer}
$
\begin{array}[t]{lll}
|A \cup B \cup C\cup D| & = & |A| + |B| + |C|+|D|\\ & & 
-|A \cap B| - |A \cap C |  -|A\cap D| - |B \cap C| -|B\cap D| - |C\cap D|\\ & & 
+|A \cap B \cap C|
+|A \cap B \cap D|
+|A \cap C \cap D|
+|B \cap C \cap D|\\ & & 
-|A \cap B \cap C\cap D|
\end{array}
$
\end{requanswer}
\end{requ}
